\chapter{Hypochondria}
\index{Hypochondria}

\section*{Definition and Overview}
Hypochondria, now formally called illness anxiety disorder, is a mental health condition in which a person is persistently worried about having or developing a serious illness, despite having no or only mild symptoms and despite reassurance that they are healthy. The concern is not relieved by negative test results and can cause significant distress and disruption to daily life. Illness anxiety disorder is a recognised and treatable condition: cognitive behavioural therapy is effective, and medications help when anxiety or depression coexists. It is different from malingering, because the worry is genuine.

\section*{Epidemiology}
Illness anxiety disorder affects around 1--5\% of the population, and it can begin at any age, often after a serious illness in the person or a family member. It is equally common in men and women. Many people with the condition see multiple doctors and undergo repeated tests, and it is associated with significant health care costs, disability, and reduced quality of life. With treatment, most people improve substantially.

\section*{Aetiology and Pathophysiology}
\begin{itemize}
    \item \textbf{Cognitive factors:} misinterpretation of normal bodily sensations as signs of disease
    \item \textbf{Attention:} excessive focus on bodily sensations amplifies them
    \item \textbf{Beliefs:} catastrophic thinking about health and illness
    \item \textbf{Life events:} past illness, death of a loved one, or childhood illness
    \item \textbf{Pathophysiology:} anxiety and attention amplify normal sensations, which reinforces the fear
    \item \textbf{Reassurance paradox:} negative test results provide only temporary relief, so the worry returns
    \item \textbf{Coexisting conditions:} anxiety, depression, and obsessive-compulsive traits are common
\end{itemize}

\section*{Clinical Features}
\begin{itemize}
    \item Persistent worry about having a serious illness
    \item Fear that persists despite normal test results
    \item Repeated checking of the body for signs of disease
    \item Frequent internet searching of symptoms
    \item Repeated visits to doctors and requests for tests
    \item Avoidance of medical information or, conversely, health care settings
    \item Significant distress and impairment in daily life
    \item Physical symptoms such as palpitations and tension from anxiety
\end{itemize}

\section*{Differential Diagnosis}
\begin{itemize}
    \item Real medical illness, which must be excluded
    \item Somatic symptom disorder, where physical symptoms dominate
    \item Panic disorder and generalised anxiety
    \item Depression, in which health worries can occur
    \item Obsessive-compulsive disorder, with intrusive health thoughts
    \item Delusional disorder, where the belief is fixed
\end{itemize}

\section*{When to Seek Medical Care}
See a GP for a structured assessment and referral to psychology or psychiatry. It is important to have a single GP coordinating care to avoid repeated unnecessary testing.

\textbf{Contact your local emergency services for:}
\begin{itemize}
    \item Thoughts of self-harm or suicide
    \item Severe panic with chest pain or collapse
    \item Any genuine emergency symptom
    \item Mental health crises
\end{itemize}

\section*{Diagnosis and Investigations}
\begin{itemize}
    \item \textbf{Clinical assessment:} a careful history of the health worries and their impact
    \item \textbf{Physical examination and tests:} to exclude genuine illness, done once and coordinated
    \item \textbf{Mental health assessment:} screening for anxiety, depression, and other conditions
    \item \textbf{Impact assessment:} effect on work, relationships, and quality of life
    \item \textbf{Longitudinal view:} the pattern of worry over time
\end{itemize}

\section*{Management}
\begin{itemize}
    \item \textbf{Cognitive behavioural therapy:} the first-line treatment, challenging health beliefs and reducing checking
    \item \textbf{Exposure therapy:} gradually facing health-related fears
    \item \textbf{Attention training:} reducing focus on bodily sensations
    \item \textbf{Medication:} SSRIs help when anxiety or depression is significant
    \item \textbf{Single-doctor care:} coordinated care reduces repeated testing
    \item \textbf{Psychoeducation:} understanding the condition reduces shame
    \item \textbf{Stress management:} relaxation and mindfulness
\end{itemize}

\section*{Complications}
\begin{itemize}
    \item Significant anxiety and distress
    \item Impaired work, relationships, and quality of life
    \item Excessive health care use and cost
    \item Unnecessary tests and procedures with their own risks
    \item Depression and social isolation
    \item Delayed treatment of genuine illness because of repeated false alarms
\end{itemize}

\section*{Prognosis and Outcomes}
Illness anxiety disorder is treatable, and most people improve with cognitive behavioural therapy. The condition may be chronic and fluctuate with stress, but with consistent treatment, health worries diminish and daily life improves. Treatment of coexisting anxiety and depression improves outcomes. Early recognition and coordinated care prevent the cycle of repeated testing that maintains the condition.

\section*{Prevention and Self-Care}
\begin{itemize}
    \item See one GP for all health concerns
    \item Ask about the value of tests before having them
    \item Limit symptom-searching on the internet
    \item Practise relaxation and mindfulness
    \item Engage in CBT and complete treatment
    \item Distract from health checking with activities
    \item Treat underlying anxiety and depression
    \item Talk openly about health worries with your doctor
\end{itemize}

\section*{Sources and Further Reading}
The following reputable sources provide further information for healthcare students and patients seeking reliable, current advice about illness anxiety disorder.

\begin{itemize}
    \item Healthdirect. \textit{Illness anxiety disorder (hypochondria).} https://www.healthdirect.gov.au/illness-anxiety-disorder
    \item Beyond Blue. \textit{Anxiety information and support.} https://www.beyondblue.org.au/
    \item Black Dog Institute. \textit{Anxiety and health worries.} https://www.blackdoginstitute.org.au/
    \item Australian Psychological Society. \textit{Finding a psychologist.} https://psychology.org.au/
    \item NHS. \textit{Health anxiety.} https://www.nhs.uk/mental-health/conditions/health-anxiety/
    \item Mayo Clinic. \textit{Illness anxiety disorder.} https://www.mayoclinic.org/
\end{itemize}
